\chapter{Process Intelligence: A Full-Lifecycle Autonomic Framework}
\label{ch:process-intelligence}

This chapter encodes the \emph{Process Intelligence} thesis as a formal
framework. The universal business of a process-intelligence system is:
\textbf{manufacturing board-admissible process reality from public standards,
type law, receipts, replay, refusal, residuals, and lifecycle authority}.
Every section that follows is a formal unpacking of that sentence.

% -----------------------------------------------------------------------------
\section{The Prime Thesis}
\label{sec:pi-thesis}
% -----------------------------------------------------------------------------

A process-intelligence system is not a reporting system that describes what
happened. It is an \emph{actuation system} that manufactures evidence-bearing
transitions across the full process lifecycle: design, simulation, construction,
activation, operation, monitoring, repair, optimization, board projection,
integration, decommissioning, and archive.

The thesis is stated as follows.

\begin{definition}[Universal Process Business]
\label{def:universal-business}
The universal business of a process-intelligence system is to manufacture
\emph{board-admissible process reality} --- that is, claims about process
behavior that are backed by public-standard evidence, type law, receipts,
replay bundles, explicit refusals, and explicit residuals --- from an
operational world $O$ by way of admissibility, actuation, projection, and
checkpoint functions.
\end{definition}

The full lifecycle addressed by the framework is:

\begin{equation}
\label{eq:lifecycle-set}
L = \{\,\mathrm{Design},\; \mathrm{Simulation},\; \mathrm{Construction},\;
      \mathrm{Activation},\; \mathrm{Operation},\; \mathrm{Monitoring},\;
      \mathrm{Repair},\; \mathrm{Optimization},\; \mathrm{BoardProjection},\;
      \mathrm{Integration},\; \mathrm{Decommission},\; \mathrm{Archive}\,\}.
\end{equation}

No lifecycle stage is privileged. Board projection is not the goal; archive is
not the end. The framework treats every stage as an equal manufacturing surface
with identical receipt and refusal obligations.

% -----------------------------------------------------------------------------
\section{Fundamental Objects}
\label{sec:pi-objects}
% -----------------------------------------------------------------------------

\begin{definition}[Fundamental Symbols]
\label{def:symbols}
The following symbols are used throughout this chapter.
\begin{itemize}
  \item $O$ --- the raw operational world (observed events, objects, relations
    as they arrive before any standard mapping).
  \item $O^*$ --- the public-standard admissible operational world: the image
    of $O$ under the admissibility and manufacturing function $\mu$.
  \item $P$ --- a process (a typed, lifecycle-carrying entity).
  \item $L$ --- the lifecycle state space (equation~\ref{eq:lifecycle-set}).
  \item $E$ --- process evidence: receipted, typed observations of process
    behavior.
  \item $T$ --- type law: the compile-time constraint system enforced by
    \texttt{wasm4pm-compat}.
  \item $R$ --- receipts: non-forgeable proofs that a lawful transition
    occurred.
  \item $F$ --- refusals: typed, law-named rejection records.
  \item $X$ --- residuals: partial, unresolved, or disputed evidence that is
    neither admitted nor refused but must be accounted for explicitly.
  \item $Q$ --- queries: structured interrogations of admitted evidence.
  \item $M$ --- mined models: process models discovered from event logs by a
    process-mining algorithm.
  \item $C$ --- conformance verdicts: outputs of model-vs-log comparison.
  \item $B$ --- board claims: projections of admitted process evidence into
    board-level representations.
  \item $D$ --- decommissioning states: terminal lifecycle positions for
    retired processes.
  \item $\mu$ --- the admissibility and manufacturing function:
    $\mu(O, T, L) \to O^*$ or $\mathrm{Refuse}(F)$.
  \item $\alpha$ --- the actuation function:
    $\alpha(K, P, L, T) \to \tau$ where $K$ is a knowledge base and $\tau$ is
    a lifecycle transition.
  \item $\rho$ --- the receipt function: $\rho(\tau) = R_\tau$, the receipt
    that proves transition $\tau$ was lawfully executed.
  \item $\pi$ --- the projection function:
    $\pi(P_i) \to B$ or $\pi(P_i, \mathrm{Risk}, \ldots) \to B_{M\&A}$.
  \item $\kappa$ --- the checkpoint function:
    $\kappa(\tau) \in \{\mathrm{ADMIT}(R),\; \mathrm{REFUSE}(F),\;
    \mathrm{PARTIAL}(X)\}$.
  \item $\delta$ --- the decommissioning function:
    $\delta(P) \to \mathrm{Retired}(P) + \mathrm{Archive}(A) +
    \mathrm{Receipt}(R_\delta)$.
\end{itemize}
\end{definition}

\subsection{Core Manufacturing Equations}

The three core equations of the framework are:

\begin{equation}
\label{eq:manufacture}
  R \vdash P_i = \mu(O^*, T, L)
\end{equation}

Equation~\eqref{eq:manufacture} states that an admissible process instance
$P_i$ is the output of applying the manufacturing function $\mu$ to the
admissible world $O^*$ under type law $T$ and lifecycle state space $L$, and
that this output is only valid when receipted by $R$.

\begin{equation}
\label{eq:board-claim}
  R,\; \mathrm{Replay},\; \mathrm{Audit} \vdash B = \pi(P_i)
\end{equation}

Equation~\eqref{eq:board-claim} states that a board claim $B$ is the
projection of a receipted, replayed, and audited process instance $P_i$.
Projection without replay is not a board claim; it is an assertion.

\begin{equation}
\label{eq:ma-claim}
  R,\; \mathrm{Replay},\; X,\; F \vdash
  B_{M\&A} = \pi(P_i,\; \mathrm{Risk},\; \mathrm{Synergy},\;
              \mathrm{Integration},\; \mathrm{Debt})
\end{equation}

Equation~\eqref{eq:ma-claim} states the M\&A deck form: a merger or
acquisition claim must project process evidence with explicit treatment of
residuals $X$ and refusals $F$ alongside risk, synergy, integration, and
technical debt. An M\&A deck that omits $X$ and $F$ is not board-admissible.

% -----------------------------------------------------------------------------
\section{The Lifecycle Calculus}
\label{sec:pi-calculus}
% -----------------------------------------------------------------------------

\begin{definition}[Lifecycle State Machine]
\label{def:lifecycle-sm}
A valid process lifecycle is a directed state machine
\[
  P \;:\; L_0 \to L_1 \to \cdots \to L_n
\]
where:
\begin{enumerate}
  \item each $L_k \in L$ (the lifecycle state space of
    equation~\ref{eq:lifecycle-set}),
  \item every transition $\tau_k : L_k \to L_{k+1}$ is \emph{lawful} (admitted
    by $\mu$ under type law $T$), and
  \item every lawful transition emits evidence: $\rho(\tau_k) = R_k$.
\end{enumerate}
There is no lawful transition that does not produce a receipt.
\end{definition}

\begin{theorem}[Transition Receipt Law]
\label{thm:receipt-law}
For every process $P$ and every transition $\tau_i$ in the lifecycle of $P$:
\[
  \forall \tau_i \in \mathrm{lifecycle}(P),\quad
  \mathrm{lawful}(\tau_i) \;\Rightarrow\;
  \exists\, R_i \text{ such that } R_i \text{ proves } \tau_i.
\]
\end{theorem}

\begin{proof}
By Definition~\ref{def:lifecycle-sm}, the receipt function $\rho$ is defined
on every lawful transition with codomain in $R$. The set $R$ of receipts is
non-empty whenever a lawful transition occurs, because $\rho$ is total on
lawful transitions by construction. A transition for which no $R_i$ exists
is therefore not a lawful transition; the lifecycle state machine does not
advance without a receipt.
\end{proof}

\subsection{The Blue River Dam Rule}

The checkpoint function $\kappa$ is total on the set of all transitions:

\begin{theorem}[Blue River Dam Rule]
\label{thm:brd-rule}
For every transition $\tau$ in any process lifecycle:
\[
  \kappa(\tau) \;\in\;
  \{\;\mathrm{ADMIT}(R),\quad \mathrm{REFUSE}(F),\quad \mathrm{PARTIAL}(X)\;\}.
\]
There is no silent success state. There is no unnamed failure state.
\end{theorem}

\begin{proof}
The codomain of $\kappa$ is defined to be exactly
$\{\mathrm{ADMIT}(R), \mathrm{REFUSE}(F), \mathrm{PARTIAL}(X)\}$ by
Definition~\ref{def:symbols}. Any runtime outcome that cannot be placed in one
of these three classes is a defect in the system implementing $\kappa$, not a
valid lifecycle outcome. Silent success is not a member of the codomain and
therefore cannot be produced by a conforming implementation.
\end{proof}

\begin{remark}
$\mathrm{PARTIAL}(X)$ is not a degraded form of $\mathrm{ADMIT}$. It is a
first-class outcome that carries explicit residuals $X$. A board claim
derived from a $\mathrm{PARTIAL}$ transition is valid precisely because the
residuals are explicit; a claim that silently discards residuals is invalid.
\end{remark}

% -----------------------------------------------------------------------------
\section{Autonomic Knowledge Actuation}
\label{sec:pi-actuation}
% -----------------------------------------------------------------------------

Classical knowledge management systems describe processes: $K$ describes $P$.
The process-intelligence framework actuates processes: $\alpha(K, P, L, T)
\to \tau$. The distinction is fundamental.

\begin{definition}[Autonomic Actuation]
\label{def:actuation}
A transition $\tau$ in the process lifecycle is \emph{valid} if and only if:
\begin{enumerate}
  \item $\tau = \alpha(K, P, L, T)$ for some knowledge base $K$, process $P$,
    current lifecycle stage $L$, and type law $T$, and
  \item $\kappa(\tau) \in \{\mathrm{ADMIT}(R),\; \mathrm{REFUSE}(F),\;
    \mathrm{PARTIAL}(X)\}$.
\end{enumerate}
A transition that satisfies condition 1 but fails condition 2 is a defect
in the actuation function $\alpha$.
\end{definition}

\begin{proposition}[Autonomic Completeness]
\label{prop:autonomic-completeness}
A process-intelligence system is \emph{autonomically complete} if and only if
for every lifecycle stage $\ell \in L$ and every lawful input $(K, P, \ell, T)$,
the actuation function $\alpha$ produces a transition $\tau$ such that
$\kappa(\tau)$ is defined and non-silent.
\end{proposition}

The contrast with prior art is captured in the following table. Old systems
accumulate knowledge and produce reports. This system actuates knowledge and
produces receipted transitions.

\begin{center}
\begin{tabular}{lll}
\toprule
\textbf{Property} & \textbf{Classical KM} & \textbf{Autonomic PI} \\
\midrule
Knowledge role & Descriptive ($K$ describes $P$) & Actuative
  ($\alpha(K,P,L,T) \to \tau$) \\
Transition output & Narrative & $\kappa(\tau) \in \{$ADMIT, REFUSE, PARTIAL$\}$ \\
Failure mode & Silent omission & $\mathrm{REFUSE}(F)$ with named law \\
Uncertainty mode & Untracked & $\mathrm{PARTIAL}(X)$ with explicit residuals \\
Board admissibility & Asserted & Receipted via $\rho(\tau) = R$ \\
\bottomrule
\end{tabular}
\end{center}

% -----------------------------------------------------------------------------
\section{Public Standards as Manufacturing Feedstock}
\label{sec:pi-standards}
% -----------------------------------------------------------------------------

Let $\mathcal{S}$ be the set of public standards recognized by the framework:

\begin{equation}
\label{eq:standards}
\mathcal{S} = \{\,\mathrm{OCEL},\; \mathrm{XES},\; \mathrm{BPMN},\;
  \mathrm{Petri},\; \mathrm{WF\text{-}net},\; \mathrm{POWL},\;
  \mathrm{Declare},\; \mathrm{ProcessTree},\; \mathrm{DFG},\;
  \mathrm{OCPQ},\; \mathrm{OTel},\; \mathrm{PROV\text{-}O},\;
  \mathrm{SHACL}\,\}.
\end{equation}

\begin{definition}[Reverse Lock-In Principle]
\label{def:reverse-lock-in}
For a process-intelligence system grounded in $\mathcal{S}$:
\[
  \frac{\partial\,\mathrm{Value}}{\partial\,\mathrm{Standards}} > 0.
\]
Every new public standard $s \notin \mathcal{S}$ is a new manufacturing
surface, not a threat to existing surfaces. Adding $s$ to $\mathcal{S}$
strictly increases the admissible operational world $O^*$.
\end{definition}

\begin{proposition}[Standards Asymmetry]
For proprietary-vendor systems, a new standard commoditizes existing lock-in
and reduces vendor margin. For a process-intelligence system satisfying
Definition~\ref{def:reverse-lock-in}, a new standard opens a new manufacturing
surface and increases the combinatorial manufacturing space
(Section~\ref{sec:pi-combinatorial}).
\end{proposition}

This asymmetry is not accidental. It follows directly from the structure of
the admissibility function $\mu$: $\mu$ accepts as input any evidence shape
that can be mapped to a public-standard type in $\mathcal{S}$. Expanding
$\mathcal{S}$ expands the domain of $\mu$.

% -----------------------------------------------------------------------------
\section{Reverse Porter Five Algebra}
\label{sec:pi-porter}
% -----------------------------------------------------------------------------

The classical Porter Five Forces model describes industry margin as a function
of five competitive forces. The process-intelligence framework inverts this
relationship.

\begin{definition}[Porter Five Force Variables]
\label{def:porter-vars}
Let $\beta$ = buyer power, $\sigma$ = supplier power, $\upsilon$ =
substitution threat, $\varepsilon$ = threat of new entrants, $\gamma$ =
competitive rivalry.
\end{definition}

\begin{definition}[Conventional Margin Equation]
\label{def:conventional-margin}
Under the classical model, industry margin is:
\[
  \mathrm{Margin} = f(-\beta,\; -\sigma,\; -\upsilon,\; -\varepsilon,\;
  -\gamma),
\]
where each force enters negatively: every increase in any force compresses
margin.
\end{definition}

\begin{theorem}[Reverse Porter Authority]
\label{thm:reverse-porter}
For a process-intelligence system satisfying Definitions~\ref{def:universal-business}
and~\ref{def:reverse-lock-in}, every increase in a Porter force increases
the demand for validation authority:
\[
  \forall\, v \in \{\beta, \sigma, \upsilon, \varepsilon, \gamma\},\quad
  v \uparrow \;\Rightarrow\;
  \mathrm{demand\_for\_validation\_authority} \uparrow.
\]
\end{theorem}

\begin{proof}
As competitive pressure increases, counterparties (buyers, regulators, boards,
auditors) increase their demand for verifiable evidence of process claims.
A system that can produce receipted, replayed, and publicly-standard-grounded
evidence satisfies this demand directly. Conventional systems that produce
narrative claims cannot. Therefore, every force that increases competitive
pressure also increases the relative advantage of the process-intelligence
framework. The proof reduces to the observation that validation authority is
not a commodity under competitive pressure; it is a differentiator.
\end{proof}

\begin{remark}
This is the formal content of the phrase ``the moat deepens as the market
heats'': in a quiet market, receipted evidence is a convenience; in a
contested market, it is the admission ticket to the board.
\end{remark}

% -----------------------------------------------------------------------------
\section{The PM4Py Oracle Formulation}
\label{sec:pi-pm4py}
% -----------------------------------------------------------------------------

PM4Py is the reference implementation of the process-mining literature. For
each PM4Py capability $c$, the framework defines a complete capability tuple:

\begin{definition}[Capability Tuple $\Omega$]
\label{def:omega}
\[
\Omega(c) = \bigl\{\;
  \mathrm{Input}(c),\;
  \mathrm{Output}(c),\;
  \mathrm{Standard}(c),\;
  \mathrm{Algorithm}(c),\;
  \mathrm{CompatType}(c),\;
  \mathrm{WasmExec}(c),\;
  \mathrm{Refusal}(c),\;
  \mathrm{Receipt}(c),\;
  \mathrm{Replay}(c),\;
  \mathrm{Lifecycle}(c),\;
  \mathrm{BoardClaim}(c)\;
\bigr\}.
\]
\end{definition}

\begin{definition}[Blue River Ready]
\label{def:brr}
A capability $c$ is \emph{Blue River Ready} ($\mathrm{BRR}(c)$) if and only if:
\[
  \mathrm{BRR}(c) \;\Longleftrightarrow\;
  \mathrm{PM4Py}(c) \;\wedge\; \mathrm{TypeLaw}(c) \;\wedge\;
  \mathrm{ExecAuthority}(c) \;\wedge\; \mathrm{Receipt}(c) \;\wedge\;
  \mathrm{Replay}(c) \;\wedge\; \mathrm{Lifecycle}(c).
\]
\end{definition}

$\mathrm{PM4Py}(c)$ means the capability is implemented in the PM4Py library.
$\mathrm{TypeLaw}(c)$ means a type-law surface exists in
\texttt{wasm4pm-compat}. $\mathrm{ExecAuthority}(c)$ means the capability has
a designated execution authority in \texttt{wasm4pm}. $\mathrm{Receipt}(c)$,
$\mathrm{Replay}(c)$, and $\mathrm{Lifecycle}(c)$ mean the capability
participates in the receipt, replay, and lifecycle frameworks respectively.

The goal of the \texttt{wasm4pm-compat} crate is to make $\mathrm{TypeLaw}(c)$
true for every $c$ for which $\mathrm{PM4Py}(c)$ is true.

% -----------------------------------------------------------------------------
\section{Six Manufacturing Algorithms}
\label{sec:pi-algorithms}
% -----------------------------------------------------------------------------

\subsection{Algorithm 1: Paper-to-Law Manufacturing}
\label{sec:pi-alg1}

For each academic paper $p$ in the process-mining literature, define the
paper extraction tuple:
\[
  \Phi(p) = \bigl\{\;
    O_p,\; A_p,\; I_p,\; Y_p,\; H_p,\;
    \mathrm{Fail}_p,\; \mathrm{Compat}_p,\; \mathrm{Wasm}_p,\;
    \mathrm{PM4Py}_p,\; \mathrm{Fixtures}_p,\; \mathrm{Receipts}_p,\;
    \mathrm{Replay}_p,\; \mathrm{Board}_p,\;
    \mathrm{Lifecycle}_p,\; \mathrm{Decommission}_p
  \;\bigr\},
\]
where $O_p$ = object types, $A_p$ = algorithm name, $I_p$ = input types,
$Y_p$ = output types, $H_p$ = hyperparameters, $\mathrm{Fail}_p$ = known
failure modes.

\begin{definition}[Coverage Law]
\label{def:coverage-law}
Paper $p$ is \emph{covered} by the framework if and only if:
\[
  \mathrm{Covered}(p) \;\Longleftrightarrow\;
  \mathrm{Compat}_p \;\vee\; \mathrm{Wasm}_p \;\vee\;
  \mathrm{ExplicitGraduationBoundary}_p.
\]
A paper that is neither in the compat surface nor the wasm surface must have
an explicit graduation boundary annotation explaining why it belongs to neither.
\end{definition}

\begin{algorithm}[H]
\caption{Paper-to-Law Manufacturing}
\label{alg:paper-law}
\DontPrintSemicolon
\KwIn{Academic paper $p$ with content corpus $\mathcal{C}_p$.}
\KwOut{$\Phi(p)$ extraction tuple; coverage verdict
       $\mathrm{Covered}(p)$ or $\mathrm{Refused}(F_p)$.}
\BlankLine
Extract $\{O_p, A_p, I_p, Y_p, H_p\}$ from abstract and formal sections of $p$\;
Extract $\mathrm{Fail}_p$ from failure-mode and complexity sections\;
\BlankLine
\tcp{Determine manufacturing surface assignment}
\uIf{structural type in \texttt{wasm4pm-compat} can represent $I_p \to Y_p$}{
  $\mathrm{Compat}_p \leftarrow \mathtt{true}$\;
}\uElseIf{algorithm $A_p$ requires execution in \texttt{wasm4pm}}{
  $\mathrm{Wasm}_p \leftarrow \mathtt{true}$\;
}\uElse{
  \uIf{explicit graduation boundary documented}{
    $\mathrm{ExplicitGraduationBoundary}_p \leftarrow \mathtt{true}$\;
  }\Else{
    \Return $\mathrm{Refused}(F_p = \texttt{MissingGraduationBoundary})$\;
  }
}
\BlankLine
Generate $\mathrm{Fixtures}_p$ (compile-fail + compile-pass) for $\Phi(p)$\;
Emit $\mathrm{Receipts}_p$ for each generated fixture\;
\Return $\Phi(p)$ with $\mathrm{Covered}(p) = \mathtt{true}$\;
\end{algorithm}

\subsection{Algorithm 2: Admissible Process Evidence}
\label{sec:pi-alg2}

\begin{algorithm}[H]
\caption{Admissible Process Evidence}
\label{alg:admit-evidence}
\DontPrintSemicolon
\KwIn{Raw value $x : T$, witness marker $W$, type law $T_{\mathrm{law}}$.}
\KwOut{\Evidence{T}{\texttt{Admitted}}{W} or $\mathrm{Refusal}(F)$ or
       $\mathrm{Partial}(X)$.}
\BlankLine
\tcp{Layer 1: structural shape}
\If{$x$ does not conform to the structural schema of $W$}{
  \Return $\mathrm{Refusal}(\texttt{StructuralSchemaMismatch})$\;
}
\tcp{Layer 2: object identity}
\If{any object identifier in $x$ is null or malformed}{
  \Return $\mathrm{Refusal}(\texttt{MalformedObjectIdentifier})$\;
}
\tcp{Layer 3: event identity}
\If{any event identifier in $x$ is null or malformed}{
  \Return $\mathrm{Refusal}(\texttt{MalformedEventIdentifier})$\;
}
\tcp{Layer 4: object-event links}
\If{any event has no qualifying object link}{
  \Return $\mathrm{Refusal}(\texttt{DanglingEventObjectLink})$\;
}
\tcp{Layer 5: object-object links}
\If{any O2O link references an undefined object}{
  \Return $\mathrm{Refusal}(\texttt{UndefinedObjectInLink})$\;
}
\tcp{Layer 6: lifecycle state}
\If{$x$ is not admissible from the current lifecycle stage $L$}{
  \Return $\mathrm{Refusal}(\texttt{InvalidLifecycleTransition})$\;
}
\tcp{Layer 7: public standard mapping}
\If{$W \notin \mathcal{S}$ and no explicit graduation boundary}{
  \Return $\mathrm{Refusal}(\texttt{UnmappedStandard})$\;
}
\tcp{Layer 8: loss policy}
\uIf{loss is present and policy is \texttt{RefuseLoss}}{
  \Return $\mathrm{Refusal}(\texttt{LossPolicyViolation})$\;
}\ElseIf{loss is present and policy is \texttt{AllowNamedProjection}
         but no name given}{
  \Return $\mathrm{Partial}(X = \{\texttt{UnnamedProjection}\})$\;
}
\BlankLine
\Return \Evidence{T}{\texttt{Admitted}}{W}$(x)$\;
\end{algorithm}

\subsection{Algorithm 3: Receipt-Bearing Execution}
\label{sec:pi-alg3}

\begin{algorithm}[H]
\caption{Receipt-Bearing Execution}
\label{alg:receipt-exec}
\DontPrintSemicolon
\KwIn{Admitted evidence \Evidence{T}{\texttt{Admitted}}{W}, algorithm $a$,
      parameters $\theta$.}
\KwOut{Verdict $V$ + Receipt $R$ + ReplayBundle $\Gamma$ + Residual $X$.}
\BlankLine
$V \leftarrow a(\Evidence{T}{\texttt{Admitted}}{W},\; \theta)$\;
\BlankLine
\tcp{Receipt law: every verdict must produce a receipt}
$R \leftarrow \rho(a,\; \theta,\; \Evidence{T}{\texttt{Admitted}}{W},\; V)$\;
\If{$R = \emptyset$}{
  \Return $\mathrm{Defect}(\texttt{MissingReceipt})$\;
  \tcp{Execution without receipt is a system defect, not a process outcome}
}
\BlankLine
\tcp{Replay bundle: receipt must be replayable against admitted evidence}
$\Gamma \leftarrow \textsc{BuildReplayBundle}(R,\;
  \Evidence{T}{\texttt{Admitted}}{W})$\;
\If{$\textsc{Replay}(\Gamma) \neq V$}{
  \Return $\mathrm{Defect}(\texttt{ReplayMismatch})$\;
}
\BlankLine
\tcp{Residual accounting: unexplained evidence must be explicit}
$X \leftarrow \Evidence{T}{\texttt{Admitted}}{W} \setminus
  \mathrm{dom}(\Gamma)$\;
\BlankLine
\tcp{Invariants: $V \Rightarrow R$, $R \Rightarrow \Gamma$,
     $\Gamma \Rightarrow \Evidence{T}{\texttt{Admitted}}{W}$}
\Return $(V,\; R,\; \Gamma,\; X)$\;
\end{algorithm}

\begin{theorem}[Execution Invariant Chain]
\label{thm:exec-invariant}
For any execution satisfying Algorithm~\ref{alg:receipt-exec}:
\[
  V \;\Rightarrow\; R, \qquad
  R \;\Rightarrow\; \Gamma, \qquad
  \Gamma \;\Rightarrow\; \Evidence{T}{\texttt{Admitted}}{W}.
\]
\end{theorem}

\begin{proof}
The algorithm returns a defect and halts if $R = \emptyset$ (no receipt from
a verdict) or if $\Gamma$ does not reproduce $V$ (replay mismatch). Therefore
any non-defect return satisfies $V \Rightarrow R$ (the receipt was produced)
and $R \Rightarrow \Gamma$ (the replay bundle was built from the receipt).
The third implication $\Gamma \Rightarrow \Evidence{T}{\texttt{Admitted}}{W}$
holds because $\Gamma$ is constructed from the admitted evidence directly;
the replay cannot reference evidence that was not admitted.
\end{proof}

\subsection{Algorithm 4: M\&A Deck Manufacturing}
\label{sec:pi-alg4}

\begin{algorithm}[H]
\caption{M\&A Deck Manufacturing}
\label{alg:ma-deck}
\DontPrintSemicolon
\KwIn{Process instance $P_i$; risk model $\mathrm{Risk}$; synergy model
      $\mathrm{Synergy}$; integration plan $\mathrm{Integration}$; debt
      registry $\mathrm{Debt}$.}
\KwOut{Board claim $B_{M\&A}$ or $\mathrm{PARTIAL}(X)$ or $\mathrm{REFUSED}(F)$.}
\BlankLine
$\mathrm{Claims} \leftarrow \textsc{ExtractClaims}(P_i,\; \mathrm{Risk},\;
  \mathrm{Synergy},\; \mathrm{Integration},\; \mathrm{Debt})$\;
\BlankLine
\For{each claim $B_j \in \mathrm{Claims}$}{
  \tcp{A claim is valid iff all six pillars hold}
  $\mathrm{valid}(B_j) \leftarrow
    E_j \wedge T_j \wedge R_j \wedge
    \mathrm{Replay}_j \wedge S_j \wedge L_j$\;
  \BlankLine
  \uIf{$\neg E_j$}{
    $B_j \leftarrow \mathrm{REFUSED}(\texttt{MissingEvidence})$\;
  }\uElseIf{$\neg T_j$}{
    $B_j \leftarrow \mathrm{REFUSED}(\texttt{TypeLawViolation})$\;
  }\uElseIf{$\neg R_j$}{
    $B_j \leftarrow \mathrm{REFUSED}(\texttt{MissingReceipt})$\;
  }\uElseIf{$\neg \mathrm{Replay}_j$}{
    $B_j \leftarrow \mathrm{REFUSED}(\texttt{ReplayUnavailable})$\;
  }\uElseIf{$\neg S_j$}{
    $B_j \leftarrow \mathrm{PARTIAL}(\texttt{UnmappedStandard})$\;
  }\uElseIf{$\neg L_j$}{
    $B_j \leftarrow \mathrm{PARTIAL}(\texttt{IncompleteLifecycle})$\;
  }
}
\BlankLine
\uIf{all claims valid}{
  \Return $B_{M\&A} = \pi(P_i,\; \mathrm{Risk},\; \mathrm{Synergy},\;
    \mathrm{Integration},\; \mathrm{Debt})$\;
}\uElseIf{some claims $\mathrm{PARTIAL}$, none $\mathrm{REFUSED}$}{
  \Return $\mathrm{PARTIAL}(X = \{B_j : B_j = \mathrm{PARTIAL}(\cdot)\})$\;
}\Else{
  \Return $\mathrm{REFUSED}(F = \{B_j : B_j = \mathrm{REFUSED}(\cdot)\})$\;
}
\end{algorithm}

\subsection{Algorithm 5: Decommissioning}
\label{sec:pi-alg5}

\begin{algorithm}[H]
\caption{Process Decommissioning}
\label{alg:decommission}
\DontPrintSemicolon
\KwIn{Active process $P$ with dependency set $\mathrm{Deps}(P)$ and
      claims set $\mathrm{Claims}(P)$.}
\KwOut{$\delta(P) = \mathrm{Retired}(P) + \mathrm{Archive}(A) +
       \mathrm{Receipt}(R_\delta)$ or $\mathrm{REFUSED}(F)$.}
\BlankLine
\tcp{Precondition 1: all dependencies closed or refused}
\ForEach{$d \in \mathrm{Deps}(P)$}{
  \If{$\kappa(d) \notin \{\mathrm{ADMIT}(\cdot),\; \mathrm{REFUSE}(\cdot)\}$}{
    \Return $\mathrm{REFUSED}(\texttt{OpenDependency}(d))$\;
  }
}
\BlankLine
\tcp{Precondition 2: all claims archived or revoked}
\ForEach{$b \in \mathrm{Claims}(P)$}{
  \If{$b$ is neither archived nor revoked}{
    \Return $\mathrm{REFUSED}(\texttt{UnarchivedClaim}(b))$\;
  }
}
\BlankLine
\tcp{Decommissioning is a lawful lifecycle transition --- it must emit a receipt}
$A \leftarrow \textsc{Archive}(P)$\;
$R_\delta \leftarrow \rho(\delta,\; P,\; A)$\;
\If{$R_\delta = \emptyset$}{
  \Return $\mathrm{REFUSED}(\texttt{MissingDecommissionReceipt})$\;
}
\BlankLine
\Return $\mathrm{Retired}(P) + \mathrm{Archive}(A) +
  \mathrm{Receipt}(R_\delta)$\;
\end{algorithm}

\subsection{Algorithm 6: Board Claim Reliability}
\label{sec:pi-alg6}

\begin{algorithm}[H]
\caption{Board Claim Reliability Assessment}
\label{alg:board-reliability}
\DontPrintSemicolon
\KwIn{Board claim $B$; evidence $E$; type law $T$; receipt $R$;
      replay bundle $\Gamma$; standard $S$; lifecycle $L$;
      residuals $X$; refusals $F$.}
\KwOut{$\mathrm{Reliable}(B)$ or $\mathrm{Unreliable}(B, \mathrm{reason})$.}
\BlankLine
\tcp{Reliability requires all eight pillars, including explicit residuals
     and known refusal boundaries}
$\mathrm{pillars} \leftarrow
  E \wedge T \wedge R \wedge \Gamma \wedge S \wedge L \wedge
  \mathrm{explicit}(X) \wedge \mathrm{explicit}(F)$\;
\BlankLine
\uIf{$\mathrm{pillars} = \mathtt{true}$}{
  \Return $\mathrm{Reliable}(B)$\;
}\Else{
  $\mathrm{missing} \leftarrow \{p : p \in \mathrm{pillars},\;
    p = \mathtt{false}\}$\;
  \Return $\mathrm{Unreliable}(B,\; \mathrm{missing})$\;
}
\BlankLine
\tcp{Note: reliability requires explicit(X), not zero(X).}
\tcp{A claim with known residuals is MORE reliable than one with hidden residuals.}
\end{algorithm}

% -----------------------------------------------------------------------------
\section{The Algebra of Board Reliance}
\label{sec:pi-algebra}
% -----------------------------------------------------------------------------

\begin{definition}[Validation Semiring]
\label{def:semiring}
Let $\mathcal{V} = \{E, T, R, \Gamma, S, L, \mathrm{explicit}(X),
\mathrm{explicit}(F)\}$ be the set of eight validation components. Define:
\begin{align*}
  v_1 \wedge v_2 &= \text{joint requirement (meet): both must hold,} \\
  v_1 \vee v_2   &= \text{sufficient condition (join): either suffices.}
\end{align*}
$(\mathcal{V}, \wedge, \vee)$ forms a semiring with
additive identity $\mathbf{0}$ (no component holds) and multiplicative
identity $\mathbf{1}$ (all components hold).
\end{definition}

\begin{theorem}[Board Reliability Condition]
\label{thm:reliability}
A board claim $B$ is reliable if and only if:
\begin{equation}
\label{eq:reliability}
  \mathrm{Reliable}(B) \;\Longleftrightarrow\;
  E \;\wedge\; T \;\wedge\; R \;\wedge\; \Gamma \;\wedge\; S \;\wedge\; L
  \;\wedge\; \mathrm{explicit}(X) \;\wedge\; \mathrm{explicit}(F).
\end{equation}
\end{theorem}

\begin{proof}
Any claim lacking evidence $E$ is ungrounded. Any claim lacking type law $T$
has no structural guarantee. Any claim lacking receipt $R$ cannot be audited.
Any claim lacking replay $\Gamma$ cannot be verified ex post. Any claim
lacking standard grounding $S$ is privately defined. Any claim lacking
lifecycle coverage $L$ may be stale. Any claim with hidden residuals ($\neg
\mathrm{explicit}(X)$) conceals uncertainty, which is worse than revealing it.
Any claim with unknown refusal boundaries ($\neg \mathrm{explicit}(F)$)
conceals the scope of what was not tested. The conjunction is therefore
necessary: removing any single conjunct creates a known reliability gap.
\end{proof}

\begin{corollary}[Residuals Strengthen Reliability]
\label{cor:residuals}
Unknown risk $>$ refused false claim. Explicitly stating residuals $X$ and
refusals $F$ does not weaken a board claim; it strengthens the epistemic
foundation of the claim by making the boundary of knowledge visible.
\end{corollary}

% -----------------------------------------------------------------------------
\section{Combinatorial Maximalism}
\label{sec:pi-combinatorial}
% -----------------------------------------------------------------------------

The manufacturing surface of the process-intelligence framework is the
Cartesian product of all admissible dimensions:

\begin{equation}
\label{eq:manufacturing-surface}
  |\mathrm{Surface}| = |\mathcal{S}| \times |L| \times |\mathcal{A}| \times
  |\mathcal{F}| \times |\mathcal{R}| \times |\mathcal{Q}| \times
  |\mathcal{B}| \times |\mathcal{D}|,
\end{equation}

where $\mathcal{S}$ is the standard set (equation~\ref{eq:standards}),
$L$ is the lifecycle state space (equation~\ref{eq:lifecycle-set}),
$\mathcal{A}$ is the algorithm set (all PM4Py capabilities), $\mathcal{F}$
is the refusal reason type set, $\mathcal{R}$ is the receipt type set,
$\mathcal{Q}$ is the query type set, $\mathcal{B}$ is the board claim type
set, and $\mathcal{D}$ is the decommissioning state set.

Each cell in this product space is a distinct manufacturing task; each
manufacturing task produces at least one receipt-bearing commit.

\subsection{Takt Equation}

\begin{definition}[Takt Equation]
\label{def:takt}
The commit throughput equation is:
\[
  C = n \times d \times t,
\]
where $C$ = total commits per month, $n$ = number of production cells (parallel
manufacturing threads), $d$ = working days per month, $t$ = commits per cell
per day.
\end{definition}

For a target of $C = 5000$ commits per month with $d = 30$ days and
$n = 10$ production cells:

\[
  t = \frac{C}{n \times d} = \frac{5000}{10 \times 30} \approx 16.7
  \text{ commits per cell per day.}
\]

This is a manufacturing rate target, not a velocity aspiration. Each commit
must carry a receipt; a commit without a receipt does not count toward $C$.
The Takt equation therefore also specifies the minimum receipt density of the
production system.

\begin{proposition}[Surface Scalability]
\label{prop:surface-scalability}
As $|\mathcal{S}|$ grows (new public standards are published), the
manufacturing surface (equation~\ref{eq:manufacturing-surface}) grows
multiplicatively. The Takt equation is unchanged; only $n$ must scale to
maintain $t$.
\end{proposition}

% -----------------------------------------------------------------------------
\section{The Blue River Dam Operating Equation}
\label{sec:pi-brd}
% -----------------------------------------------------------------------------

The complete operating cycle of the process-intelligence framework is
expressed as a function composition:

\begin{equation}
\label{eq:brd}
  \mathrm{BlueRiverDam} = \kappa \circ \rho \circ \alpha \circ \mu.
\end{equation}

Expanding equation~\eqref{eq:brd}:

\begin{equation}
\label{eq:brd-expanded}
  \kappa\bigl(\rho\bigl(\alpha\bigl(\mu(O^*)\bigr)\bigr)\bigr)
  \;\to\; \mathrm{ALIVE} \;\mid\; \mathrm{PARTIAL} \;\mid\;
  \mathrm{REFUSED}.
\end{equation}

Reading left to right: raw operational world $O$ is passed through $\mu$ to
produce the admissible operational world $O^*$; the actuation function
$\alpha$ produces a lifecycle transition from $O^*$; the receipt function
$\rho$ stamps the transition with a non-forgeable proof; the checkpoint
function $\kappa$ classifies the result into one of the three and only three
terminal states.

\begin{theorem}[Operating Equation Completeness]
\label{thm:brd-completeness}
The Blue River Dam operating equation~\eqref{eq:brd-expanded} is complete:
every possible outcome of a process-intelligence manufacturing cycle is
captured by $\{\mathrm{ALIVE}, \mathrm{PARTIAL}, \mathrm{REFUSED}\}$.
\end{theorem}

\begin{proof}
By the Blue River Dam Rule (Theorem~\ref{thm:brd-rule}), $\kappa$ maps every
transition to exactly one of $\{\mathrm{ADMIT}(R), \mathrm{REFUSE}(F),
\mathrm{PARTIAL}(X)\}$. $\mathrm{ADMIT}(R)$ is the receipted success case,
written $\mathrm{ALIVE}$ at the system level. $\mathrm{REFUSE}(F)$ is the
named refusal case, written $\mathrm{REFUSED}$. $\mathrm{PARTIAL}(X)$ is the
explicit residual case. The three cases are mutually exclusive and exhaustive
by the definition of $\kappa$'s codomain.
\end{proof}

\begin{remark}
The four-function composition $\kappa \circ \rho \circ \alpha \circ \mu$ is
not an architectural diagram. It is a mathematical specification. An
implementation that produces output in $\{\mathrm{ALIVE}, \mathrm{PARTIAL},
\mathrm{REFUSED}\}$ for every input in $O$ satisfies the specification. An
implementation that can produce output outside this set --- including null,
exception, or silent success --- does not satisfy the specification and is a
defect.
\end{remark}

% -----------------------------------------------------------------------------
\section{Chapter Summary}
\label{sec:pi-summary}
% -----------------------------------------------------------------------------

This chapter has established the following formal content.

\paragraph{Lifecycle calculus.} A valid process lifecycle is a directed state
machine over $L$ (equation~\ref{eq:lifecycle-set}) in which every lawful
transition emits a receipt (Theorem~\ref{thm:receipt-law}). There is no silent
success state (Theorem~\ref{thm:brd-rule}).

\paragraph{Autonomic actuation.} The actuation function $\alpha(K,P,L,T) \to
\tau$ replaces the classical knowledge-management relationship $K$ describes
$P$ with a receipt-bearing transition obligation (Definition~\ref{def:actuation}).

\paragraph{Reverse Porter authority.} Every increase in competitive pressure
increases the demand for validation authority; the process-intelligence
framework satisfies this demand through receipted, replayed,
publicly-standard-grounded evidence (Theorem~\ref{thm:reverse-porter}).

\paragraph{PM4Py oracle.} Blue River Ready is a six-predicate conjunction
over PM4Py coverage, type law, execution authority, receipt, replay, and
lifecycle (Definition~\ref{def:brr}).

\paragraph{Six algorithms.} Paper-to-law manufacturing
(Algorithm~\ref{alg:paper-law}), admissible evidence
(Algorithm~\ref{alg:admit-evidence}), receipt-bearing execution
(Algorithm~\ref{alg:receipt-exec}), M\&A deck manufacturing
(Algorithm~\ref{alg:ma-deck}), decommissioning
(Algorithm~\ref{alg:decommission}), and board claim reliability
(Algorithm~\ref{alg:board-reliability}) together constitute the complete
manufacturing protocol.

\paragraph{Board reliability algebra.} Reliability is a conjunction of eight
components (Theorem~\ref{thm:reliability}); explicit residuals strengthen
rather than weaken reliability (Corollary~\ref{cor:residuals}).

\paragraph{Operating equation.} $\mathrm{BlueRiverDam} = \kappa \circ \rho
\circ \alpha \circ \mu$ is complete and exhaustive
(Theorem~\ref{thm:brd-completeness}).

\bigskip
\noindent\textbf{The product is CodeManufactory; process reality is merely
proof that CodeManufactory works.}
